\documentclass[10pt,a4paper]{amsart}
\usepackage[margin=19mm]{geometry}
\usepackage{amsmath,amssymb,mathtools}
\usepackage[scr=rsfs,cal=euler]{mathalfa}
\usepackage{xfrac}
\usepackage[unicode,hidelinks,bookmarks=false]{hyperref}
\newcommand{\ie}{i.e.\ }
\newcommand{\cf}{cf.\ }
\newcommand{\resp}{respectively\ }
\newcommand{\Subsection}{\subsection}
\providecommand{\nobreakdash}{\nobreak\hbox{-}}
\setlength{\emergencystretch}{2em}
\multlinegap0pt
\usepackage{bm}
\usepackage{bbm}
\usepackage{dsfont}
\usepackage{xspace}
\usepackage{cancel}
\usepackage{mathtools}
\usepackage{amsthm}
\makeatletter
\@ifundefined{theorem}{\newtheorem{theorem}{Theorem}}{}
\@ifundefined{lemma}{\newtheorem{lemma}[theorem]{Lemma}}{}
\makeatother
\newcommand{\E}{\mathbb{E}}
\newcommand{\cE}{\mathscr{E}}
\title[Moments of sums of exponentials, beyond CHS]{Moments of sums of exponentials, beyond CHS}
\author{Silouanos Brazitikos, Colin Tang, Tomasz Tkocz}
\date{Source: arXiv:2602.03058v1 (CC BY 4.0)}
\begin{document}
\maketitle

\noindent\emph{Source and licence.} Moments of sums of exponentials, beyond CHS. Version arXiv:2602.03058v1 (CC BY 4.0), distributed under the Creative Commons Attribution 4.0 International licence, as recorded in the arXiv licence metadata for this exact version and shown on the abstract page https://arxiv.org/abs/2602.03058v1. Full rights notice: licenses/arxiv-2602.03058-CC-BY-4.0.txt.\par\smallskip

\noindent\textit{Editorial scope.} Counted unit: the Lemma whose source label is \texttt{lm:F} in the article (source0.tex lines 605--607, source label \texttt{lm:F}), delivered together with the article's own complete original proof of it (source0.tex lines 614--682, one \texttt{proof} environment of 69 lines). No mathematics is written, repaired or re-derived: statement and proof are the article's own text line for line. The proof is detached from the statement in the source: the article states the result at lines 605--607 and prints its proof at lines 614--682, and the proof environment's own header names the statement, so the association is the article's own. Required context retained uncounted: none. Every definition, display and label of those lines is retained unchanged. Omitted from this package and not redistributed: 1--604, 608--613, 683--1000. Nothing inside a retained excerpt is omitted or altered: every retained body is delivered exactly as the article's lines are written, so no omission receipt of any kind accompanies this package. Citations: the retained excerpts carry no citation command at all, so no bibliography entry of the article is transcribed and no reference is redistributed. Reference adaptation: none was needed and none was made; every reference inside the delivered unit resolves inside it, so no unresolved reference or citation is produced. Printed slips kept literal: no printed word, symbol, hypothesis or number of the counted statement or of its proof was altered. Layout and class: the article's own class is \texttt{article} with packages including inputenc, amssymb, amsthm, amsmath, mathrsfs, enumerate, graphicx, xcolor, setspace, hyperref, mathalpha; the wrapper is the lane's pinned \texttt{amsart} editorial wrapper at 10pt on A4, which restates the theorem-like environments the retained text needs and adds the article's own macro definitions that the retained text needs (\texttt{\textbackslash E}, \texttt{\textbackslash cE}), with extra wrapper packages (bm, bbm, dsfont, xspace, cancel, mathtools). The delivered numbering is the unit's own. Assets: no retained span contains a figure environment, a graphics-inclusion command, a PDF-page inclusion, a table or a TikZ picture; no image file is copied into this package, the package declares no asset dependency and no image file is redistributed with it. Licence: version arXiv:2602.03058v1 of this article is distributed under the Creative Commons Attribution 4.0 International licence, as recorded in the arXiv licence metadata for this exact version and shown on the abstract page https://arxiv.org/abs/2602.03058v1. A full rights notice is delivered at licenses/arxiv-2602.03058-CC-BY-4.0.txt. Characters: every retained excerpt is pure ASCII, so the delivered engine, XeLaTeX with its default Latin Modern text font, reports no missing character.
\begin{lemma}\label{lm:F}
Functions $F_\ell$ are Schur-convex for $\ell = 0, 1, 2, 3$.
\end{lemma}

\begin{proof}[Proof of Lemma \ref{lm:F}]
We have, $\E \cE_j^m = m!$, thus
\begin{align*}
P(x) &= \E \exp\left(-\sum \sqrt{x_j}\cE_j\right) = \prod \frac{1}{1+\sqrt{x_j}}, \\
M_1(x) &= \E \left(\sum \sqrt{x_j}\cE_j\right) = \sum \sqrt{x_j}, \\
M_2(x) &= \E \left(\sum \sqrt{x_j}\cE_j\right)^2 = 2\sum x_j + \sum_{j \neq k} \sqrt{x_jx_k}, \\
M_3(x) &= \E \left(\sum \sqrt{x_j}\cE_j\right)^3 = 6\sum x_j^{3/2} + 6\sum_{j \neq k} x_j\sqrt{x_k} + \sum_{i \neq j \neq k} \sqrt{x_ix_jx_k}.
\end{align*}
We shall use the Ostrowski criterion, so we will also need the partial derivatives,
\begin{align*}
\frac{\partial P}{\partial x_1} &= \frac{-\frac{1}{2\sqrt{x_1}}}{(1+\sqrt{x_1})^2}\prod_{j\neq 1} \frac{1}{1+\sqrt{x_j}} = -\frac{1}{2\sqrt{x_1}(1+\sqrt{x_1})}P(x),\\
\frac{\partial M_1}{\partial x_1} &= \frac{1}{2\sqrt{x_1}},\\
\frac{\partial M_2}{\partial x_1} &= 2 + \frac{1}{\sqrt{x_1}}\sum_{j\neq 1}\sqrt{x_j} = 1+\frac{M_1(x)}{\sqrt{x_1}},\\
\frac{\partial M_3}{\partial x_1} &= 9\sqrt{x_1} + 6\sum_{j\neq 1}\sqrt{x_j} + \frac{3}{\sqrt{x_1}}\sum_{j\neq 1}x_j + \frac{3}{2\sqrt{x_1}}\sum_{j\neq k \neq 1} \sqrt{x_jx_k} \\
&= 3M_1(x) + \frac{3}{\sqrt{x_1}}\sum x_j + \frac{3}{2\sqrt{x_1}}\sum_{j\neq k \neq 1} \sqrt{x_jx_k}\\
&=3\sqrt{x_1}+\frac{3}{\sqrt{x_1}}\sum x_j + \frac{3}{2\sqrt{x_1}}\sum_{j\neq k} \sqrt{x_jx_k}
\end{align*}
For convenience, we will also use the variables $b_j \equiv \sqrt{x_j}$.

\emph{Case $\ell = 0$.} We have, $Q_0(t) = 1 - e^{-t}$, thus
\[
F_0(x) = 1 - P(x),
\]
therefore,
\[
\frac{\partial F_0}{\partial x_1} - \frac{\partial F_0}{\partial x_2} = \frac12\left(\frac{1}{b_1(1+b_1)} - \frac{1}{b_2(1+b_2)}\right)P(x) = \frac{(b_2-b_1)(1+b_1+b_2)}{2b_1b_2(1+b_1)(1+b_2)}P(x).
\]
When $x_1 > x_2$, equivalently $b_1 > b_2$, this difference of partial derivatives is negative, so $F_0$ is Schur-concave.

\emph{Case $\ell = 1$.}
We have, $Q_1(t) = e^{-t} -1 + t$, thus
\[
F_1(x) = P(x)-1 + M_1(x),
\]
therefore,
\begin{align*}
\frac{\partial F_1}{\partial x_1} - \frac{\partial F_1}{\partial x_2} &= -\frac{(b_2-b_1)(1+b_1+b_2)}{2b_1b_2(1+b_1)(1+b_2)}P(x) + \frac{b_2-b_1}{2b_1b_2} \\
&=\left(\frac{b_2-b_1}{2b_1b_2}\right)\left(1 - \frac{1+b_1+b_2}{(1+b_1)(1+b_2)}P(x)\right).
\end{align*}
Plainly $\frac{1+b_1+b_2}{(1+b_1)(1+b_2)} < 1$ and $P(x) < 1$, so the second bracket is positive.
As a result, when $x_1 > x_2$, equivalently $b_1 > b_2$, this difference of partial derivatives is negative, so $F_1$ is Schur-concave.


\emph{Case $\ell = 2$.}
We have, $Q_2(t) = -e^{-t} +1 - t + \frac{t^2}{2}$, thus
\[
F_2(x) = -P(x) + 1 - M_1(x) + \frac{1}{2}M_2(x),
\]
therefore,
\begin{align*}
\frac{\partial F_2}{\partial x_1} - \frac{\partial F_2}{\partial x_2} &= \frac{(b_2-b_1)(1+b_1+b_2)}{2b_1b_2(1+b_1)(1+b_2)}P(x) - \frac{b_2-b_1}{2b_1b_2} + \frac{b_2-b_1}{2b_1b_2}M_1(x)  \\
&=\left(\frac{b_2-b_1}{2b_1b_2}\right)\left(\frac{1+b_1+b_2}{(1+b_1)(1+b_2)}P(x)-1 + M_1(x)\right).
\end{align*}
This is perhaps less obvious, but the second bracket is positive.

\textbf{Claim.} For positive numbers $b_1, \dots, b_n$, we have
\[
\frac{1+b_1+b_2}{(1+b_1)(1+b_2)} > \left(1 - \sum b_j\right)\prod (1+b_j).
\]
As a result, when $x_1 > x_2$, equivalently $b_1 > b_2$, this difference of partial derivatives is negative, so $F_2$ is Schur-concave.

\emph{Case $\ell = 3$.} Again, we calculate that
\begin{align*}
	&\frac{\partial F_3}{\partial x_1} - \frac{\partial F_3}{\partial x_2}=\\ &\left(\frac{b_2-b_1}{2b_1b_2}\right)\left(6\sum b_i^2+3\sum_{i\neq j}b_ib_j-6b_1b_2-\frac{1+b_1+b_2}{(1+b_1)(1+b_2)}P(x)+1-M_1(x)\right).
\end{align*}
Note that $$6\sum b_i^2+3\sum_{i\neq j}b_ib_j-6b_1b_2-M_1(x)=3+3\left(\sum b_i\right)^2-6b_1b_2-\sum b_i\geq 0,$$
since $$6b_1b_2\leq\frac{3}{2}(b_1+b_2)^2\leq \frac{3}{2}(\sum b_i)^2$$ and 
$$\frac{1}{2}(\sum b_i)^2+\frac{1}{2}\geq\sum b_i.$$
\end{proof}

\end{document}
